% =============================================================================
% Complete Presentation Script
% Additive Channels in Curved Fitness Landscapes
% 31 Slides with Full Speaker Notes
% Estimated Duration: 50–60 Minutes
%
% Daniel Ortiz-Barrientos & Mark Cooper
% December 2024
% =============================================================================

\documentclass[aspectratio=169,11pt]{beamer}

% -----------------------------------------------------------------------------
% Packages
% -----------------------------------------------------------------------------
\usepackage{amsmath,amssymb,amsthm}
\usepackage{graphicx}
\usepackage{booktabs}
\usepackage{xcolor}
\usepackage{tikz}
\usetikzlibrary{arrows.meta,positioning,shapes,calc}
\usepackage{pgfplots}
\pgfplotsset{compat=1.17}
\usepackage{hyperref}
\usepackage{pgfpages}

% -----------------------------------------------------------------------------
% Speaker Notes Configuration
% Toggle between presentation mode and notes mode
% -----------------------------------------------------------------------------
% Uncomment ONE of the following lines:
%\setbeameroption{hide notes}                        % Presentation mode (no notes)
\setbeameroption{show notes on second screen=right}  % Dual screen: slides left, notes right
%\setbeameroption{show only notes}                   % Notes only (for script review)

% -----------------------------------------------------------------------------
% Theme and Colors
% -----------------------------------------------------------------------------
\usetheme{default}
\usecolortheme{default}

% Custom colors
\definecolor{UQpurple}{RGB}{81,36,122}
\definecolor{channelgreen}{RGB}{76,153,0}
\definecolor{escapeamber}{RGB}{255,153,0}
\definecolor{curveblue}{RGB}{51,102,204}

\setbeamercolor{title}{fg=UQpurple}
\setbeamercolor{frametitle}{fg=UQpurple}
\setbeamercolor{structure}{fg=UQpurple}
\setbeamercolor{block title}{bg=UQpurple,fg=white}
\setbeamercolor{block body}{bg=UQpurple!10}

% -----------------------------------------------------------------------------
% Math Commands
% -----------------------------------------------------------------------------
\newcommand{\var}{\mathrm{Var}}
\newcommand{\cov}{\mathrm{Cov}}
\newcommand{\tr}{\mathrm{tr}}
\newcommand{\Ag}{\mathcal{A}_g}
\newcommand{\Vlin}{V_{\mathrm{lin}}}
\newcommand{\Vquad}{V_{\mathrm{quad}}}
\newcommand{\defeq}{\coloneqq}

% -----------------------------------------------------------------------------
% Title Information
% -----------------------------------------------------------------------------
\title[Additive Channels]{Additive Channels in Curved Fitness Landscapes}
\subtitle{Why do additive models work so well when genomes are so complicated?}
\author{Daniel Ortiz-Barrientos \and Mark Cooper}
\institute{The University of Queensland}
\date{December 2024}

% =============================================================================
\begin{document}
% =============================================================================

% -----------------------------------------------------------------------------
% TITLE SLIDE
% -----------------------------------------------------------------------------
\begin{frame}
\titlepage
\end{frame}

\note[itemize]{
\item \textbf{[1 minute]} Good morning, everyone. Thank you for coming.
\item Today I want to talk about a puzzle that has been sitting at the heart of evolutionary genetics for decades---a puzzle that I think we can now resolve.
\item The puzzle is this: \textbf{Why do additive models work so well, when we know that genomes are incredibly complicated?}
\item Think about it. For over a century, we've been using models that assume genetic effects simply add up. The breeder's equation. GBLUP. Genomic selection. Fisher's fundamental theorem. These tools assume the world is linear---and they work. They predict evolutionary change. They improve crops and livestock. They've been validated in thousands of studies.
\item And yet, every time we look closely at the molecular details, we find epistasis. Gene interactions. Genotype-by-environment effects. Diminishing returns. The genome is clearly not additive.
\item So how can both things be true? That's what we're going to figure out today.
}

% =============================================================================
% PART I: THE PARADOX (Slides 1–6)
% =============================================================================
\section{Part I: The Paradox}

% -----------------------------------------------------------------------------
% SLIDE 2: The Empirical Success of Additive Models
% -----------------------------------------------------------------------------
\begin{frame}{The Empirical Success of Additive Models}
\frametitle{The Empirical Success of Additive Models}

\textbf{What We Observe in Practice}

\begin{itemize}
    \item Breeder's equation works: $\Delta\bar{z} = h^2 S$
    \item GBLUP predicts breeding values with high accuracy
    \item Responses to selection follow $\mathbf{G}\boldsymbol{\beta}$
    \item Heritabilities are stable across generations
\end{itemize}

\vspace{0.5cm}
\textit{Visual: Plot of predicted vs realised selection response across species/traits.}

\end{frame}

\note[itemize]{
\item \textbf{[1.5 minutes]} Let's start with what we actually observe in practice.
\item This figure shows predicted versus realised response to selection across a wide range of studies---maize, cattle, wheat, Drosophila, mice. The points cluster remarkably close to the one-to-one line. When we predict how much a population should change using an additive model, we get it right most of the time.
\item The breeder's equation---that simple formula relating response to heritability and selection differential---has been delivering practical results since the 1930s. Dairy cattle produce twice as much milk as they did fifty years ago. Crop yields have increased steadily. We've been selecting on complex traits for generations, and the additive machinery keeps working.
\item GBLUP, genomic selection, all of these modern tools---they're built on additive assumptions. And they achieve prediction accuracies of 0.6, 0.7, sometimes higher. That's not noise. That's real predictive power.
\item So the first truth we need to acknowledge is that \textbf{additive models work}. Whatever we say about molecular complexity, we can't deny that linear predictions are empirically successful.
}

% -----------------------------------------------------------------------------
% SLIDE 3: The Molecular Reality
% -----------------------------------------------------------------------------
\begin{frame}{The Molecular Reality}
\frametitle{The Molecular Reality}

\textbf{What We Know About Genetic Architecture}

\begin{itemize}
    \item Epistasis is pervasive in fitness landscapes
    \item G$\times$E interaction is ubiquitous
    \item Diminishing returns: mutations have smaller effects on fitter backgrounds
    \item Sign epistasis: effects can reverse in different contexts
\end{itemize}

\vspace{0.5cm}
\textit{Visual: Examples from deep mutational scanning, microbial evolution, or crop G$\times$E studies.}

\end{frame}

\note[itemize]{
\item \textbf{[1.5 minutes]} Now let's look at the other side of the puzzle.
\item Panel A shows diminishing returns epistasis---this is data from microbial evolution experiments. As you fix more beneficial mutations, each new mutation has a smaller effect. The fitness gain from your fifth mutation is smaller than from your first. That's nonlinear. That's epistasis.
\item Panel B shows genotype-by-environment interaction. These reaction norms cross. The best genotype in environment one is not the best in environment two. Rankings change. If effects were simply additive, these lines would be parallel.
\item Panel C shows pairwise epistasis from a fitness landscape study. The off-diagonal elements are not zero. The effect of mutation A depends on whether mutation B is present.
\item And this is not cherry-picked data. Deep mutational scanning, experimental evolution, crop G$\times$E studies---everywhere we look with molecular resolution, we find interactions.
\item \textbf{The genome is demonstrably non-additive.}
\item So we have two truths that seem to contradict each other.
}

% -----------------------------------------------------------------------------
% SLIDE 4: The Paradox Stated
% -----------------------------------------------------------------------------
\begin{frame}{The Paradox Stated}
\frametitle{The Paradox Stated}

\textbf{Two Apparently Contradictory Truths}

\begin{enumerate}
    \item Additive models predict evolutionary/breeding outcomes with remarkable accuracy
    \item The underlying genetic architecture is demonstrably non-additive
\end{enumerate}

\vspace{0.8cm}
\begin{block}{The Question}
Under what conditions should populations behave as if the world were additive, and how do they leave that regime?
\end{block}

\end{frame}

\note[itemize]{
\item \textbf{[1 minute]} So here's the paradox, stated plainly.
\item On the left: what we observe. The breeder's equation works. GBLUP predicts accurately. Heritabilities are stable. Fisher's theorem holds. Check, check, check, check.
\item On the right: what molecular biology shows us. Epistasis is pervasive. G$\times$E is ubiquitous. Diminishing returns. Sign epistasis. X, X, X, X.
\item How can both be true?
\item This is the central question of our paper: \textbf{Under what conditions should populations behave as if the world were additive? And how do they leave that regime?}
\item Notice I'm not asking whether genomes are additive---we know they're not. I'm asking when populations \emph{behave} additively despite the underlying complexity. That's a different question, and it has a geometric answer.
}

% -----------------------------------------------------------------------------
% SLIDE 5: Preview of the Answer
% -----------------------------------------------------------------------------
\begin{frame}{Preview of the Answer}
\frametitle{Preview of the Answer}

\textbf{Additivity is Local, Not Global}

\begin{itemize}
    \item A curved landscape can look linear if you only sample a small region
    \item Populations with narrow genetic variance stay in ``additive channels''
    \item The key is the ratio of linear to quadratic variance contributions
    \item Escape from additivity is a geometric event, not a mysterious breakdown
\end{itemize}

\vspace{0.5cm}
\textit{Visual: Curved surface with a small ellipse (population) where tangent plane matches.}

\end{frame}

\note[itemize]{
\item \textbf{[1.5 minutes]} Here's the preview of our answer.
\item Look at this figure. The curve is the true fitness landscape---it's clearly nonlinear. But see that small ellipse near the centre? That's a population with limited genetic variance. And see how well the tangent line matches the curve in that region?
\item If you're a population sitting in that green zone, the world looks linear to you. You don't ``see'' the curvature because you're not sampling enough of the landscape to detect it.
\item This is the key insight: \textbf{additivity is local, not global}. It's not a property of the genome. It's a property of how much of the fitness landscape your population samples.
\item Think of it like driving on a curved road. If you only drive a hundred metres at a time, every road feels straight. You don't notice the curvature. Only when you try to drive a long distance does the curve become apparent.
\item Populations with narrow genetic variance are like drivers going short distances. They stay in what we call ``additive channels''---regions where the linear approximation is good enough.
\item The question becomes: when do you leave the channel? And what happens when you do?
}

% -----------------------------------------------------------------------------
% SLIDE 6: Talk Outline
% -----------------------------------------------------------------------------
\begin{frame}{Talk Outline}
\frametitle{Talk Outline}

\textbf{What We'll Cover}

\begin{enumerate}
    \item The framework: Taylor expansion of the fitness landscape
    \item The additivity index: A diagnostic for when linear models work
    \item Additive channels: Regions where Fisher's theorem holds
    \item Escape dynamics: How populations leave additive regimes
    \item Applications: Breeding programs and adaptive radiations
    \item Empirical implications: What to measure and when to worry
\end{enumerate}

\end{frame}

\note[itemize]{
\item \textbf{[30 seconds]} Here's where we're going.
\item First, I'll show you the mathematical framework---how a Taylor expansion of the fitness landscape connects to classical quantitative genetics.
\item Then I'll introduce the additivity index---a number between zero and one that tells you how ``linear'' the world looks to your population.
\item We'll define additive channels and show that Fisher's theorem and the breeder's equation are local properties of these channels.
\item Then we'll explore escape dynamics---how populations leave additive channels, and why this matters for breeding and evolution.
\item Finally, I'll discuss applications to breeding programs and adaptive radiations, and give you practical diagnostics for your own data.
\item Let's start with the framework.
}

% =============================================================================
% PART II: THE FRAMEWORK (Slides 7–12)
% =============================================================================
\section{Part II: The Framework}

% -----------------------------------------------------------------------------
% SLIDE 7: The Genotype–Phenotype–Fitness Map
% -----------------------------------------------------------------------------
\begin{frame}{The Genotype--Phenotype--Fitness Map}
\frametitle{The Genotype--Phenotype--Fitness Map}

\textbf{Setting Up the Problem}

Three spaces connected by two maps:
\[
\text{Genotype } \mathbf{g} \xrightarrow{\phi} \text{Phenotype } \mathbf{z} \xrightarrow{\psi} \text{Fitness } W
\]

Combined map:
\[
F(\mathbf{g}, e) = \psi(\phi(\mathbf{g}, e), e)
\]

\vspace{0.5cm}
\textit{Visual: Diagram showing genotype cloud $\to$ phenotype cloud $\to$ fitness values.}

\end{frame}

\note[itemize]{
\item \textbf{[1.5 minutes]} Let me set up the mathematical framework.
\item We have three spaces. Genotype space---this is where allele frequencies and genetic values live. Phenotype space---the traits we measure. And fitness---the single number that determines evolutionary fate.
\item Two maps connect these spaces. The genotype-phenotype map, which we call $\phi$---this is development, physiology, everything that turns DNA into an organism. And the fitness function $\psi$---this is natural selection, the relationship between what you look like and how many offspring you leave.
\item The combined map $F$ takes you directly from genotype to fitness. And crucially, the environment $e$ enters both maps. Development depends on environment. Fitness depends on environment. This is where G$\times$E lives.
\item Now, here's the key move. We're not going to assume $F$ is linear. We're going to approximate it locally and ask: when is the linear approximation good enough?
\item That's what the Taylor expansion does.
}

% -----------------------------------------------------------------------------
% SLIDE 8: Taylor Expansion – The Core Idea
% -----------------------------------------------------------------------------
\begin{frame}{Taylor Expansion -- The Core Idea}
\frametitle{Taylor Expansion -- The Core Idea}

\textbf{Any Smooth Function is Locally Polynomial}

\[
F(\mathbf{g}) \approx \underbrace{F_0}_{\text{baseline}} + \underbrace{\mathbf{b}^\top \Delta\mathbf{g}}_{\text{linear (additive)}} + \underbrace{\frac{1}{2}\Delta\mathbf{g}^\top \mathbf{H} \Delta\mathbf{g}}_{\text{quadratic (epistasis)}} + \cdots
\]

\begin{itemize}
    \item $\mathbf{b} = \nabla F$ is the gradient (additive effects)
    \item $\mathbf{H}$ is the Hessian (curvature/epistasis)
    \item Close to reference: linear dominates. Far away: curvature matters.
\end{itemize}

\end{frame}

\note[itemize]{
\item \textbf{[2 minutes]} The Taylor expansion is the mathematical engine of everything that follows.
\item The idea is simple: any smooth function can be approximated by polynomials. Near any point, we can write fitness as a baseline, plus a linear term, plus a quadratic term, plus higher-order terms we'll ignore.
\item Look at the left panel. The black curve is the true fitness function---clearly nonlinear. The gray line is just the value at the reference point---useless. The blue line is the linear approximation---it captures the slope. The green curve is the quadratic approximation---it captures the curvature too.
\item Near the reference point, all three approximations agree. But as you move away, they diverge. The linear approximation drifts. The quadratic tracks the true function longer.
\item Now here's the biological interpretation. The gradient $\mathbf{b}$---these are your additive effects. This is what breeders estimate. The Hessian $\mathbf{H}$---this is curvature. The diagonal elements relate to dominance-like effects. The off-diagonal elements are epistasis.
\item So the Taylor expansion doesn't assume linearity. It \emph{decomposes} the landscape into linear and nonlinear parts. Then we can ask: which part matters more for my population?
}

% -----------------------------------------------------------------------------
% SLIDE 9: Connecting to Quantitative Genetics
% -----------------------------------------------------------------------------
\begin{frame}{Connecting to Quantitative Genetics}
\frametitle{Connecting to Quantitative Genetics}

\textbf{Classical Quantities Emerge from the Expansion}

\begin{center}
\begin{tabular}{ll}
\toprule
\textbf{Taylor term} & \textbf{QG interpretation} \\
\midrule
$\mathbf{b}_g = \mathbf{J}_g^\top \mathbf{a}_z$ & Additive genetic effects \\
$\mathbf{H}_{gg} = \mathbf{J}_g^\top \mathbf{H}_{zz} \mathbf{J}_g$ & Epistasis (genetic curvature) \\
$\mathbf{H}_{ge}$ & Genotype-by-environment interaction \\
$\mathbf{b}_g^\top \boldsymbol{\Sigma}_g \mathbf{b}_g$ & Additive genetic variance \\
\bottomrule
\end{tabular}
\end{center}

\vspace{0.5cm}
\begin{block}{}
The breeder's equation is the linear term of this expansion!
\end{block}

\end{frame}

\note[itemize]{
\item \textbf{[1.5 minutes]} Here's the beautiful part: classical quantitative genetics falls out of this framework automatically.
\item The gradient in genotype space, $\mathbf{b}_g$, is the product of the Jacobian of the genotype-phenotype map and the selection gradient on phenotypes. These are your additive genetic effects---exactly what breeders estimate as allelic effects or breeding values.
\item The genetic Hessian $\mathbf{H}_{gg}$ captures how the fitness surface curves in genotype space. This is epistasis---how the effect of one allele depends on the genetic background.
\item The mixed Hessian $\mathbf{H}_{ge}$ is where G$\times$E lives. It captures how additive effects change across environments.
\item And when you combine the gradient with the genetic covariance matrix $\boldsymbol{\Sigma}_g$, you get the additive genetic variance in fitness: $\mathbf{b}_g^\top \boldsymbol{\Sigma}_g \mathbf{b}_g$. That's $V_A$.
\item The breeder's equation---$\Delta\bar{z} = \mathbf{G}\boldsymbol{\beta}$---is just the linear term of this expansion. It works when the linear term dominates. It fails when the quadratic term matters.
\item So classical QG isn't wrong---it's \emph{local}. It's the first-order approximation to something more complex.
}

% -----------------------------------------------------------------------------
% SLIDE 10: From Derivatives to Variance Components
% -----------------------------------------------------------------------------
\begin{frame}{From Derivatives to Variance Components}
\frametitle{From Derivatives to Variance Components}

\textbf{How Population Variation Creates Fitness Variation}

Under Gaussian assumptions:
\begin{align}
\Vlin &= \mathbf{b}^\top \boldsymbol{\Sigma}_g \mathbf{b} & \text{(additive variance)} \\[0.5em]
\Vquad &= \frac{1}{2}\tr[(\mathbf{H}\boldsymbol{\Sigma}_g)^2] & \text{(epistatic variance)}
\end{align}

\vspace{0.5cm}
\begin{alertblock}{Key insight}
$\Vlin \propto \sigma^2$ but $\Vquad \propto \sigma^4$
\end{alertblock}

\end{frame}

\note[itemize]{
\item \textbf{[2 minutes]} Now we connect derivatives to variance components. This is where the key insight emerges.
\item The Taylor expansion tells us how fitness depends on genotype for an individual. But quantitative genetics asks about populations---about variation. So we need to combine the derivatives with the covariance structure of standing genetic variation.
\item Under Gaussian assumptions---which we need for tractable math---the variance of the linear term is $\mathbf{b}^\top \boldsymbol{\Sigma}_g \mathbf{b}$. This is just a quadratic form. It scales as $\sigma^2$---the variance of the genetic distribution.
\item The variance of the quadratic term involves the trace of a product of matrices. The details are in the appendix, but the key point is this: it scales as $\sigma^4$. That's the \emph{square} of the variance.
\item Look at the figure on the left. When genetic variance is small, the green bar dominates. $\Vlin$ is most of the fitness variance. But as genetic variance grows, the red bar catches up. $\Vquad$ grows as $\sigma^4$, which is faster than $\sigma^2$.
\item \textbf{This is the scaling insight that resolves the paradox.} When populations have narrow genetic variance, the quadratic term is negligible because $\sigma^4$ is tiny. The world looks linear. As variance expands, curvature becomes detectable.
}

% -----------------------------------------------------------------------------
% SLIDE 11: The Additivity Index
% -----------------------------------------------------------------------------
\begin{frame}{The Additivity Index}
\frametitle{The Additivity Index}

\textbf{A Diagnostic for ``How Linear Is the World?''}

\[
\Ag = \frac{\Vlin}{\Vlin + \Vquad}
\]

\begin{itemize}
    \item $\Ag \approx 1$: Linear terms dominate; additive models work
    \item $\Ag \approx 0$: Curvature dominates; additive models fail
    \item $\Ag$ is \textbf{local}: it depends on where the population sits and how spread out it is
\end{itemize}

\end{frame}

\note[itemize]{
\item \textbf{[1.5 minutes]} This brings us to the central diagnostic: the additivity index.
\item It's simply the fraction of fitness variance explained by the linear term. $\Vlin$ over $\Vlin$ plus $\Vquad$.
\item When $\Ag$ is close to one, almost all fitness variance comes from the linear term. The quadratic contribution is negligible. In that regime, additive models work. The breeder's equation predicts accurately. Fisher's theorem holds.
\item When $\Ag$ is close to zero, curvature dominates. The quadratic term generates most of the fitness variance. Additive predictions fail. Rankings change unexpectedly. Epistasis matters.
\item Look at the contour plot on the left. Near the centre, $\Ag$ is high---green. As you move away from the reference point, $\Ag$ drops---red. The additive approximation gets worse.
\item But here's the crucial point: $\Ag$ is not a property of the landscape alone. It depends on where the population sits and how spread out it is. The same landscape can have high $\Ag$ for a narrow population and low $\Ag$ for a wide one.
\item This is why additivity is local, not global.
}

% -----------------------------------------------------------------------------
% SLIDE 12: Ag as R²
% -----------------------------------------------------------------------------
\begin{frame}{$\Ag$ as $R^2$}
\frametitle{$\Ag$ as $R^2$}

\textbf{The Additivity Index is the $R^2$ of the Linear Predictor}

Under the Gaussian assumptions:
\[
\Ag = \text{Corr}(L, \tilde{\ell})^2 = R^2
\]

where $L = \mathbf{b}^\top \Delta\mathbf{g}$ is the linear predictor and $\tilde{\ell} = L + Q$ is true log-fitness.

\vspace{0.5cm}
\begin{block}{Interpretation}
$\Ag$ tells you what fraction of fitness variation your additive model explains.
\end{block}

\end{frame}

\note[itemize]{
\item \textbf{[1.5 minutes]} Here's a result that makes $\Ag$ directly interpretable.
\item Under the Gaussian assumptions, $\Ag$ equals the squared correlation between the linear predictor and true fitness. It's an $R^2$. It's the prediction accuracy of your additive model, squared.
\item Look at the scatter plots. On the left, high $\Ag$: the linear predictor $L$ tracks true fitness closely. Points cluster around the one-to-one line. $R^2 = 0.91$.
\item On the right, low $\Ag$: the scatter is huge. The linear predictor explains little of the fitness variation. $R^2 = 0.37$.
\item The geometric interpretation is in panel C. Think of fitness as a vector that decomposes into a linear component $L$ and a quadratic component $Q$. These are orthogonal under Gaussianity. $\Ag$ is the cosine squared of the angle between total fitness and its linear projection.
\item This connects directly to GBLUP. The prediction accuracy of GBLUP is bounded above by $\sqrt{\Ag}$. If $\Ag$ is 0.81, your GBLUP accuracy can't exceed 0.9, no matter how much data you have. This tells you whether your bottleneck is estimation error or model misspecification.
}

% =============================================================================
% PART III: ADDITIVE CHANNELS (Slides 13–18)
% =============================================================================
\section{Part III: Additive Channels}

% -----------------------------------------------------------------------------
% SLIDE 13: Defining Additive Channels
% -----------------------------------------------------------------------------
\begin{frame}{Defining Additive Channels}
\frametitle{Defining Additive Channels}

\textbf{Regions Where Additive Models Work}

An additive channel is a connected region of genotype space where:
\begin{enumerate}
    \item $\Ag$ exceeds a threshold (e.g., 0.8 or 0.9)
    \item Response to selection follows $\mathbf{G}\boldsymbol{\beta}$
    \item Directional heritability is high in the direction of selection
\end{enumerate}

\vspace{0.5cm}
\textit{Channels are corridors carved by curvature through which populations can evolve ``additively.''}

\end{frame}

\note[itemize]{
\item \textbf{[1.5 minutes]} Now we can define what we mean by an additive channel.
\item An additive channel is a connected region of genotype space where the additivity index exceeds some threshold---say 0.8 or 0.9. Within the channel, responses to selection follow the breeder's equation. Directional heritability predicts response. Fisher's theorem holds.
\item Look at this 3D visualization. The blue surface is the fitness landscape. The green tube is the additive channel---a corridor where the linear approximation is good. The amber ellipse is a population sitting comfortably within the channel.
\item The metaphor is deliberate. Channels are carved by the curvature of the landscape. Where curvature is mild, the channel is wide. Where curvature is strong, the channel narrows or bends. Populations can cruise through channels using simple additive machinery. But if they stray outside, they need to account for the curvature.
\item The threshold is a modelling choice. Setting $\Ag^0 = 0.9$ demands that 90\% of fitness variance come from the linear term. That's stringent. $\Ag^0 = 0.8$ is more permissive. The right threshold depends on your application and your tolerance for prediction error.
}

% -----------------------------------------------------------------------------
% SLIDE 14: Fisher's Theorem as a Corollary
% -----------------------------------------------------------------------------
\begin{frame}{Fisher's Theorem as a Corollary}
\frametitle{Fisher's Theorem as a Corollary}

\textbf{The Fundamental Theorem Holds Within Additive Channels}

\[
\text{Fisher's FTNS:} \quad \frac{d\bar{W}}{dt} = V_A(W)
\]

This follows from our framework when:
\begin{itemize}
    \item $\Ag \approx 1$ (curvature contribution is small)
    \item We consider only the additive genetic component
\end{itemize}

\vspace{0.5cm}
\begin{alertblock}{}
Fisher's theorem is not a universal law; it's a local property of high-$\Ag$ regions.
\end{alertblock}

\end{frame}

\note[itemize]{
\item \textbf{[1.5 minutes]} This brings us to a conceptual shift in how we think about classical results.
\item Fisher's fundamental theorem---``the rate of increase in mean fitness equals the additive genetic variance in fitness''---is often treated as a deep truth about adaptation. A law of nature.
\item But in our framework, it's a corollary. It follows from the Taylor expansion when the additivity index is high. When $\Ag \approx 1$, the quadratic term contributes little to fitness variance. The Price equation reduces to Fisher's result.
\item Look at the right panel. I've simulated a population evolving within an additive channel. The blue dots show the observed rate of fitness increase. The green line shows $V_A(W)$, the prediction from Fisher's theorem. They track each other beautifully.
\item But this only works because $\Ag$ is high throughout the trajectory. If the population left the channel, if curvature started to matter, Fisher's prediction would fail.
\item So Fisher's theorem is not universal. It's local. It's a property of additive channels, not of evolution in general. This doesn't diminish Fisher---it contextualizes him. It tells us when his theorem applies and when it doesn't.
}

% -----------------------------------------------------------------------------
% SLIDE 15: Directional Heritability and Response
% -----------------------------------------------------------------------------
\begin{frame}{Directional Heritability and Response}
\frametitle{Directional Heritability and Response}

\textbf{Why the Breeder's Equation Works (Sometimes)}

Within an additive channel:
\[
\Delta\bar{z} = \mathbf{G}\boldsymbol{\beta}
\]

Directional heritability $h^2_d(\mathbf{u}) = \frac{\mathbf{u}^\top \mathbf{G} \mathbf{u}}{\mathbf{u}^\top \mathbf{P} \mathbf{u}}$ predicts response.

\vspace{0.5cm}
\textbf{When $\Ag$ drops:}
\begin{itemize}
    \item Observed responses deviate from predictions
    \item Apparent changes in $\mathbf{G}$ over short timescales
    \item Epistasis becomes visible in rankings
\end{itemize}

\end{frame}

\note[itemize]{
\item \textbf{[1.5 minutes]} The breeder's equation tells us that response equals the $\mathbf{G}$ matrix times the selection gradient. Directional heritability tells us how well the population can respond in any particular direction.
\item Within an additive channel, these tools work. If you know $\mathbf{G}$ and $\boldsymbol{\beta}$, you can predict response. If selection pushes in a direction with high $h^2_d$, you get rapid change. If it pushes in a constrained direction, response is slow.
\item But when $\Ag$ drops, these predictions break down. You see the symptoms on the right side of the slide.
\item Observed responses deviate systematically from $\mathbf{G}\boldsymbol{\beta}$. You predict one thing, you observe another. Breeders call this ``unexpected re-ranking.''
\item The $\mathbf{G}$ matrix appears to change over short timescales. You estimate $\mathbf{G}$ in generation one, estimate it again in generation five, and the eigenvectors have rotated. But this isn't necessarily evolution changing the genetic architecture---it might be the population moving through regions of the landscape where the local structure looks different.
\item Epistasis becomes visible. Genotypes that ranked well in one context rank poorly in another. The additivity assumption is failing.
\item These are warning signs that you're leaving an additive channel.
}

% -----------------------------------------------------------------------------
% SLIDE 16: The 1D Explicit Condition
% -----------------------------------------------------------------------------
\begin{frame}{The 1D Explicit Condition}
\frametitle{The 1D Explicit Condition}

\textbf{When Is a Population in an Additive Channel?}

For a single genetic coordinate with variance $V_A$:
\[
\Ag = \frac{1}{1 + \frac{b^2 V_A}{2a^2}}
\]

\textbf{Channel condition} ($\Ag \geq \Ag^0$):
\[
\boxed{V_A \leq \frac{2a^2}{b^2} \cdot \frac{1 - \Ag^0}{\Ag^0}}
\]

Three controls: selection strength ($a$), curvature ($b$), variance ($V_A$).

\end{frame}

\note[itemize]{
\item \textbf{[2 minutes]} Let me make this concrete with the simplest possible case: one genetic dimension.
\item Suppose fitness depends quadratically on a single trait: $W = az - \frac{1}{2}bz^2$. Here $a$ is the strength of directional selection---how steeply fitness increases with the trait. And $b$ is the curvature---how quickly that increase slows down.
\item For a population with additive genetic variance $V_A$, the additivity index works out to this simple formula: one over one plus a ratio involving $b$, $V_A$, and $a$.
\item Look at the left panel. As $V_A$ increases, $\Ag$ decreases. More genetic variance means you're sampling more of the landscape, hitting more curvature. The green and amber lines show different thresholds---where you'd say you're in or out of the channel.
\item The boxed equation is the channel condition. It says: you're in an additive channel if and only if your genetic variance is below a critical threshold. That threshold depends on how curved the landscape is ($b$), how strong directional selection is ($a$), and how stringent you want to be ($\Ag^0$).
\item Let me work through a numerical example. If $a = 1$, $b = 0.5$, and you want $\Ag > 0.9$, then $V_A$ must be less than about 0.89. If your actual $V_A$ is 2, your $\Ag$ drops to 0.8---you're at the edge of the channel.
}

% -----------------------------------------------------------------------------
% SLIDE 17: Biological Interpretation of the Condition
% -----------------------------------------------------------------------------
\begin{frame}{Biological Interpretation of the Condition}
\frametitle{Biological Interpretation of the Condition}

\textbf{What Controls Entry to/Exit from Additive Channels?}

\begin{center}
\begin{tabular}{lll}
\toprule
\textbf{Factor} & \textbf{Effect on $\Ag$} & \textbf{Biological interpretation} \\
\midrule
$\uparrow$ Curvature ($b$) & $\downarrow$ & Strongly curved landscapes are harder \\
$\uparrow$ Genetic variance ($V_A$) & $\downarrow$ & Wider sampling hits curvature \\
$\uparrow$ Directional selection ($a$) & $\uparrow$ & Slope dominates over curvature \\
\bottomrule
\end{tabular}
\end{center}

\vspace{0.5cm}
\begin{block}{}
Bottlenecked populations are more likely to be ``additive.''
\end{block}

\end{frame}

\note[itemize]{
\item \textbf{[1.5 minutes]} Let's translate the math into biological intuition.
\item Three factors control whether you're in an additive channel.
\item First, \textbf{curvature}. Strongly curved landscapes are harder to approximate linearly. If the fitness surface bends sharply, even a narrow population will feel the curvature. Think of a rugged landscape with many epistatic interactions---that's high $b$, low $\Ag$.
\item Second, \textbf{genetic variance}. This is the big one. Wide populations sample more of the landscape. They ``see'' the curvature. Narrow populations stay local. They experience only the slope. This is why bottlenecked populations behave more additively---they have less variance to expose the nonlinearity.
\item Third, \textbf{directional selection}. Strong directional selection means a steep slope. The linear term is large. Even if there's curvature, the slope dominates. Weak directional selection means the slope is small, and curvature becomes relatively more important.
\item Here's a surprising implication: populations under strong selection are more likely to behave additively, not less. The slope overwhelms the curvature. This might explain why selection experiments work so well in early generations, when response is fast and the linear term dominates.
}

% -----------------------------------------------------------------------------
% SLIDE 18: The Self-Correcting Dynamic
% -----------------------------------------------------------------------------
\begin{frame}{The Self-Correcting Dynamic}
\frametitle{The Self-Correcting Dynamic}

\textbf{Selection Pushes Populations Into Additive Channels}

The key insight: $\Vquad \propto \sigma^4$ shrinks faster than $\Vlin \propto \sigma^2$.

As selection compresses variance:
\begin{itemize}
    \item $\Vquad$ drops as $\sigma^4$
    \item $\Vlin$ drops as $\sigma^2$
    \item $\Ag$ rises
\end{itemize}

\vspace{0.5cm}
\begin{alertblock}{}
Additivity is self-reinforcing under selection!
\end{alertblock}

\end{frame}

\note[itemize]{
\item \textbf{[2 minutes]} Now here's the most optimistic result in the paper.
\item Selection doesn't just use additive channels---it \emph{creates} them. There's a self-correcting dynamic that pushes populations into additive regimes.
\item Here's why. Selection compresses genetic variance. That's the Bulmer effect, among other things. As you select, the genetic variance in your population decreases.
\item Now, $\Vlin$ scales as $\sigma^2$. $\Vquad$ scales as $\sigma^4$. When variance shrinks, both terms shrink---but the quadratic term shrinks faster. Much faster. If variance halves, $\Vlin$ drops by half, but $\Vquad$ drops by three-quarters.
\item Look at the left panel. At the start, when variance is high, there's substantial red---quadratic variance. But as breeding cycles proceed and variance compresses, the red shrinks faster than the green. By cycle 10, it's almost all green.
\item The middle panel shows $\Ag$ rising over time. The population enters the additive channel and stays there. By cycle 5, $\Ag$ exceeds 0.8. By cycle 10, it's above 0.95.
\item This explains why additive models work even on complex landscapes. Selection is self-correcting. It compresses variance into regimes where the linear approximation is valid. The population naturally finds and stays in additive channels.
\item This is, I think, the resolution of the paradox. Epistatic architectures produce additive-looking dynamics because selection acts.
\item One important caveat, though: this improves linear predictability, but it can also reduce $\Vlin$ and therefore slow realised gain if $\mathbf{G}$ or $|\mathbf{b}|$ becomes small. Compression helps additivity but can hurt response. We'll come back to this tension when we discuss breeding applications.
}

% =============================================================================
% PART IV: ESCAPE FROM ADDITIVITY (Slides 19–22)
% =============================================================================
\section{Part IV: Escape from Additivity}

% -----------------------------------------------------------------------------
% SLIDE 19: How Populations Leave Additive Channels
% -----------------------------------------------------------------------------
\begin{frame}{How Populations Leave Additive Channels}
\frametitle{How Populations Leave Additive Channels}

\textbf{The Geometry of Escape}

Rate of change of $\Ag$ along an evolutionary trajectory:
\[
\frac{d\Ag}{dt} \propto \nabla\Ag^\top \boldsymbol{\Sigma}_g \mathbf{b}
\]

\begin{itemize}
    \item \textbf{Positive}: moving deeper into the channel
    \item \textbf{Negative}: approaching escape
    \item \textbf{Zero}: moving along the channel boundary
\end{itemize}

\end{frame}

\note[itemize]{
\item \textbf{[1.5 minutes]} If selection pushes populations into additive channels, how do they ever leave?
\item The answer lies in this equation. The rate of change of $\Ag$ along an evolutionary trajectory is the dot product of the gradient of $\Ag$ with the direction of evolutionary motion.
\item Think of $\Ag$ as defining contours over genotype space---like elevation contours on a map. High $\Ag$ is the centre of the valley; low $\Ag$ is the edge. The gradient of $\Ag$ points toward higher additivity---toward the centre.
\item The evolutionary velocity---the direction the population moves---is determined by $\boldsymbol{\Sigma}_g \mathbf{b}$, the product of genetic variance and selection gradient.
\item If these point in the same direction, you move deeper into the channel. $\Ag$ increases. Additive predictions improve.
\item If they point in opposite directions, you're escaping. $\Ag$ decreases. You're crossing contours toward lower additivity.
\item Look at the left panel. The green trajectory stays in the high-$\Ag$ region. The red trajectory crosses contours and escapes into low-$\Ag$ territory.
\item Escape is a geometric event. It happens when selection drives the population across $\Ag$ contours rather than along them.
}

% -----------------------------------------------------------------------------
% SLIDE 20: Environmental Rotation
% -----------------------------------------------------------------------------
\begin{frame}{Environmental Rotation}
\frametitle{Environmental Rotation}

\textbf{G$\times$E as Gradient Rotation}

Environmental change rotates the selection gradient relative to genetic variance:
\[
\Vlin = \sum_i \lambda_i (\mathbf{u}_i^\top \mathbf{b}(e))^2
\]

When $\mathbf{b}(e)$ rotates away from leading eigenvectors of $\mathbf{G}$:
\begin{itemize}
    \item $\Vlin$ drops (selection pushes in low-variance directions)
    \item $\Ag$ falls even if curvature and total variance are unchanged
\end{itemize}

\end{frame}

\note[itemize]{
\item \textbf{[2 minutes]} One major route of escape is through environmental change.
\item This figure shows the $\mathbf{G}$ matrix as an ellipse---elongated in the direction where genetic variance is high. The red arrow is the selection gradient---the direction fitness increases.
\item In environment 1, the gradient aligns with the major axis of the ellipse. Selection pushes in a direction rich in genetic variance. $\Vlin$ is large. $\Ag$ is high: 0.92.
\item In environment 2, the gradient has rotated. It's no longer aligned with the major axis. Selection now pushes partly in directions where variance is low.
\item In environment 3, the misalignment is severe. The gradient is nearly perpendicular to the major axis. $\Vlin$ has dropped dramatically. $\Ag$ is now 0.45---well outside the additive channel.
\item And here's the key point: I haven't changed the curvature. I haven't changed the total genetic variance. I've only rotated the selection gradient. The same genetic architecture can look additive or non-additive depending on which direction you're trying to push.
\item This is G$\times$E viewed geometrically. The environment rotates the selection gradient relative to the eigenstructure of $\mathbf{G}$. When alignment is good, additivity holds. When alignment is poor, the linear term shrinks and curvature becomes relatively more important.
}

% -----------------------------------------------------------------------------
% SLIDE 21: Introgression and Variance Expansion
% -----------------------------------------------------------------------------
\begin{frame}{Introgression and Variance Expansion}
\frametitle{Introgression and Variance Expansion}

\textbf{Wide Crosses Disrupt Additive Channels}

When you introgress exotic germplasm:
\begin{itemize}
    \item $\mathbf{G}$ expands (more genetic variance)
    \item $\mathbf{G}$ may rotate (new variance directions)
    \item $\Vquad$ grows as $\sigma^4$ -- faster than $\Vlin$
    \item $\Ag$ drops sharply
\end{itemize}

\vspace{0.5cm}
\textbf{Prediction:} GBLUP accuracy drops post-introgression, then recovers as selection re-compresses variance.

\end{frame}

\note[itemize]{
\item \textbf{[2 minutes]} Another major route of escape is through introgression or wide crosses.
\item Look at the left panel. The small green ellipse is your elite germplasm---narrow variance, high $\Ag$, GBLUP works great. Now you introgress exotic material. The pink ellipse is the expanded population---much more variance, possibly in new directions.
\item What happens? $\Vquad$ scales as $\sigma^4$. When you double variance, $\Vlin$ doubles but $\Vquad$ quadruples. The ratio shifts dramatically. $\Ag$ drops.
\item The middle panel shows the trajectory. Before introgression, $\Ag$ is rising as selection compresses variance. Then you make the wide cross---boom, $\Ag$ drops from 0.83 to 0.5. You've left the additive channel.
\item But here's the good news: it recovers. Selection re-compresses variance. $\Ag$ climbs back up. By cycle 17, you're back in the channel.
\item This has direct implications for breeding programs. Wide crosses are necessary for long-term gain---you need new variation. But they temporarily break your additive models. GBLUP accuracy will drop. Rankings will change unexpectedly. You need to be patient and let selection do its work, or adjust your models.
\item The framework predicts this pattern: drop, then recovery. And it tells you why.
}

% -----------------------------------------------------------------------------
% SLIDE 22: The Driving Analogy
% -----------------------------------------------------------------------------
\begin{frame}{The Driving Analogy}
\frametitle{The Driving Analogy}

\textbf{Cruise Control vs Active Steering}

Think of additive channels as straight roads:
\begin{itemize}
    \item On a straight road, cruise control (linear models) works fine
    \item At a curve, you need active steering (nonlinear terms)
    \item Most driving (evolution) is on straight sections
    \item Curves are transient -- you negotiate them and find a new straight section
\end{itemize}

\vspace{0.5cm}
\textit{Adaptive radiations, range expansions, and breeding into new environments are the ``curves'' where quadratic terms spike.}

\end{frame}

\note[itemize]{
\item \textbf{[1.5 minutes]} Let me give you an analogy that might help consolidate these ideas.
\item Think of additive channels as straight roads. When you're on a straight road, cruise control works perfectly. You set it and forget it. The car maintains speed without any intervention. That's what additive models are. They're cruise control. They work when the road is straight.
\item At a curve, cruise control isn't enough. You need to actively steer. You need to slow down, turn the wheel, negotiate the bend. That's when you need nonlinear terms---when the landscape curves.
\item But here's the thing: most driving is on straight sections. Even on a winding mountain road, you spend most of your time on relatively straight stretches, with brief periods of active steering at the curves.
\item Evolution is similar. Populations spend most of their time in additive channels---straight roads where Fisher's theorem and the breeder's equation work. Occasionally, they encounter curves: adaptive radiations, range expansions, introgression events, novel environments. During these transients, curvature matters. Then they settle into a new straight section.
\item The curves are real. The quadratic terms spike. But they're transient. The population negotiates the curve and finds a new additive channel.
}

% =============================================================================
% PART V: APPLICATIONS (Slides 23–28)
% =============================================================================
\section{Part V: Applications}

% -----------------------------------------------------------------------------
% SLIDE 23: Applications in Plant Breeding
% -----------------------------------------------------------------------------
\begin{frame}{Applications in Plant Breeding}
\frametitle{Applications in Plant Breeding}

\textbf{$\Ag$ as a Breeding Decision Metric}

Use $\Ag$ as a dashboard indicator:

\begin{center}
\begin{tabular}{ll}
\toprule
\textbf{$\Ag$ value} & \textbf{Action} \\
\midrule
High ($> 0.85$) & Standard GBLUP; invest in data quality \\
Moderate ($0.6$--$0.85$) & Expect re-ranking; consider interactions \\
Low ($< 0.6$) & Change model, not just data \\
\bottomrule
\end{tabular}
\end{center}

\vspace{0.3cm}
$\Ag$ tells you whether your bottleneck is estimation error or model misspecification.

\end{frame}

\note[itemize]{
\item \textbf{[2 minutes]} Let me turn to practical applications, starting with plant breeding.
\item Think of $\Ag$ as a dashboard indicator, like the fuel gauge in your car. It tells you something important about the state of your system.
\item High $\Ag$---above 0.85---means you're in a comfortable additive channel. Standard GBLUP will work well. Your bottleneck is probably data quality: more plots, better phenotyping, larger training populations. Invest there.
\item Moderate $\Ag$---between 0.6 and 0.85---means you're near the edge. Expect some re-ranking across environments or cycles. You might want to consider G$\times$E models, or at least be cautious about predictions.
\item Low $\Ag$---below 0.6---means you've left the additive channel. More data won't help much because the model is misspecified. You need to change the model: add quadratic terms, fit interactions explicitly, use environment-specific models.
\item Here's the key diagnostic question: \textbf{Is my bottleneck estimation error or model misspecification?}
\item If $\Ag$ is high but accuracy is low, you need more data. The model is right but poorly estimated.
\item If $\Ag$ is low and accuracy is low, you need a better model. More data won't fix a misspecified model.
\item This is something prediction accuracy alone can't tell you. $\Ag$ adds crucial information.
\item But there's an important distinction we need to make: high $\Ag$ tells you about predictability, not about the magnitude of gains. Let me show you what I mean.
}

% -----------------------------------------------------------------------------
% SLIDE 24: Why Programs Plateau Even Inside Additive Channels
% -----------------------------------------------------------------------------
\begin{frame}{Why Programs Plateau Even Inside Additive Channels}
\frametitle{Why Programs Plateau Even Inside Additive Channels}

\textbf{High $\Ag$ = Good Predictability, Not Guaranteed Gain}

\begin{itemize}
    \item Additivity can remain high while response slows
    \item Response depends on slope and aligned variance: $\Vlin = \mathbf{b}^\top \mathbf{G} \mathbf{b}$
    \item Plateau routes:
    \begin{itemize}
        \item $|\mathbf{b}| \downarrow$ (approaching local optimum)
        \item $\mathbf{G} \downarrow$ (variance depletion, Bulmer, drift)
    \end{itemize}
\end{itemize}

\vspace{0.3cm}
\textit{Visual: Left panel: $\Ag$ rising. Right panel: Gain per cycle declining.}

\end{frame}

\note[itemize]{
\item \textbf{[1--1.5 minutes]} One important clarification for breeders: being deep in an additive channel does not guarantee indefinite progress.
\item High $\Ag$ means curvature contributes little to the region you're sampling---so linear ranking stays reliable. Your GBLUP predictions are accurate. But realised gain depends on the local slope and the genetic variance aligned with it: $\Vlin = \mathbf{b}^\top \mathbf{G} \mathbf{b}$.
\item Over cycles, two things often happen.
\item First, as the population improves, the local gradient can shrink. You're approaching a local optimum, or a ridge, or diminishing returns kick in. The slope $|\mathbf{b}|$ drops.
\item Second, selection and relatedness compress $\mathbf{G}$---sometimes dramatically. The Bulmer effect within generations, linkage disequilibrium decay, genetic drift in small populations. All of these reduce aligned variance.
\item Either effect drives $\Vlin$ down and gains slow, even while prediction remains accurate. This is why elite programs can look highly additive---rankings stable, GBLUP working well---and yet still plateau.
\item The solution? New variation through introgression, or a shifted breeding target that finds a new slope. That's how you restart long-run gain.
}

% -----------------------------------------------------------------------------
% SLIDE 25: MET Pipeline Integration
% -----------------------------------------------------------------------------
\begin{frame}{MET Pipeline Integration}
\frametitle{MET Pipeline Integration}

\textbf{Operationalising $\Ag$ in a Breeding Program}

You already estimate:
\begin{itemize}
    \item $\mathbf{G}$ from your MET mixed model
    \item $\mathbf{b}$ from your selection index or TPE weights
\end{itemize}

You add:
\begin{itemize}
    \item $\mathbf{H}$ from quadratic regression on genetic PCs
    \item Compute $\Ag$ with bootstrap uncertainty
\end{itemize}

\vspace{0.3cm}
\textbf{Annual routine:} Plot $\Ag$ alongside cross-validated accuracy. Flag drops.

\end{frame}

\note[itemize]{
\item \textbf{[1.5 minutes]} How do you actually compute $\Ag$ in a breeding program?
\item Good news: you already have most of what you need. Your MET mixed model gives you the $\mathbf{G}$ matrix---the genetic covariance among environments or traits. Your selection index or target population of environments defines the direction $\mathbf{b}$.
\item What you add is an estimate of curvature. The simplest approach: fit a quadratic regression of EBVs or phenotypes on the leading genetic principal components. The quadratic coefficients give you an approximation to the Hessian.
\item Then you compute $\Ag$ using the formulas I showed. Use bootstrap or Bayesian methods to get uncertainty intervals.
\item The annual routine is simple. Every cycle, plot $\Ag$ alongside your cross-validated accuracy. Watch for drops. If $\Ag$ drops sharply after an introgression event, don't panic---it's expected. If $\Ag$ drops gradually over time, investigate. Maybe the breeding target has shifted. Maybe you're pushing into new phenotypic space.
\item The point isn't to hit a specific $\Ag$ target. It's to have the information. To know when you're in a channel and when you're leaving.
}

% -----------------------------------------------------------------------------
% SLIDE 26: Simulation Validation
% -----------------------------------------------------------------------------
\begin{frame}{Simulation Validation}
\frametitle{Simulation Validation}

\textbf{Breeding Simulation: $\Ag$ Predicts GBLUP Accuracy}

Simulated breeding program with epistatic landscape:
\begin{itemize}
    \item Selection intensity varies: stronger selection $\to$ higher $\Ag$
    \item Introgression drops $\Ag$, then recovery
    \item Environmental blur affects accuracy at all $\Ag$ levels
\end{itemize}

\vspace{0.3cm}
\textbf{Key result:} Realised GBLUP accuracy tracks $\sqrt{\Ag}$ as upper bound.

\end{frame}

\note[itemize]{
\item \textbf{[1.5 minutes]} We've validated these ideas with breeding simulations.
\item We created an epistatic fitness landscape---multiple QTLs with pairwise interactions---and ran simulated breeding programs with varying selection intensity, introgression events, and environmental noise.
\item The left panel shows the key result. Each point is one simulation scenario. The x-axis is $\Ag$; the y-axis is realised GBLUP accuracy. The red dashed line is $\sqrt{\Ag}$---the theoretical upper bound.
\item The points cluster below the bound, as expected. Accuracy can't exceed the limit set by model specification. But the correlation is strong: $r = 0.78$. $\Ag$ predicts accuracy remarkably well.
\item The right panel breaks this down by scenario. Strong selection gives high $\Ag$ and high accuracy. Weak selection gives lower values. Post-introgression scenarios show the drop and recovery pattern we predicted.
\item Environmental blur---measurement error, G$\times$E not captured by the model---reduces accuracy at all $\Ag$ levels. But the relationship with $\Ag$ persists.
\item This validates the framework. $\Ag$ really does predict how well additive models will work.
}

% -----------------------------------------------------------------------------
% SLIDE 27: Adaptive Radiations
% -----------------------------------------------------------------------------
\begin{frame}{Adaptive Radiations}
\frametitle{Adaptive Radiations}

\textbf{Escape from Additivity in Natural Populations}

Adaptive radiations involve:
\begin{itemize}
    \item Colonisation of new environments (gradient rotation)
    \item Hybridisation (variance expansion)
    \item Exploration of phenotypic novelty (curvature matters)
\end{itemize}

\vspace{0.5cm}
\textbf{Prediction:} Early radiation phases show low $\Ag$ and poor additive prediction; later phases stabilise into new channels.

\end{frame}

\note[itemize]{
\item \textbf{[1.5 minutes]} The framework applies to natural populations too, particularly to adaptive radiations.
\item Radiations involve exactly the processes that cause escape from additive channels. Colonisation of new environments rotates the selection gradient. Hybridisation expands variance. Exploration of phenotypic novelty pushes populations into curved regions of the landscape.
\item Think of cichlid radiations in African lakes, Darwin's finches, Hawaiian silverswords. These are bursts of diversification where populations explore new ecological niches, new morphologies, new ways of life.
\item Our framework predicts that early radiation phases should show low $\Ag$. Lineages are in the ``curves''---navigating complex fitness landscapes where epistasis and G$\times$E matter. Additive predictions would fail.
\item Later phases, after lineages have settled into their niches, should show higher $\Ag$. Variance has compressed. Local adaptation has aligned the population with a new additive channel. Fisher's theorem starts working again.
\item This is testable. With time-series data on $\mathbf{G}$ matrices, selection gradients, and predictive accuracy, you could track $\Ag$ through a radiation and see if it follows this pattern.
}

% -----------------------------------------------------------------------------
% SLIDE 28: Empirical Diagnostics
% -----------------------------------------------------------------------------
\begin{frame}{Empirical Diagnostics}
\frametitle{Empirical Diagnostics}

\textbf{Detecting Additive Channels in Real Data}

Warning signs that you're leaving an additive channel:
\begin{enumerate}
    \item Prediction accuracy drops systematically
    \item Leading eigenvectors of $\mathbf{G}$ rotate rapidly
    \item Observed response deviates from $\mathbf{G}\boldsymbol{\beta}$
    \item Environmental sensitivity increases unexpectedly
    \item Realized heritability drops
\end{enumerate}

\vspace{0.3cm}
\textit{These are measurable from standard breeding/evolutionary data.}

\end{frame}

\note[itemize]{
\item \textbf{[1.5 minutes]} Let me give you a practical checklist of warning signs.
\item First, prediction accuracy drops systematically. Not random noise---a consistent decline over cycles or generations. That suggests you're drifting out of an additive channel.
\item Second, the leading eigenvectors of $\mathbf{G}$ rotate rapidly. If your genetic PCs are moving substantially between estimates, the local structure of the landscape is changing. That could mean you're traversing, not cruising.
\item Third, observed response deviates from $\mathbf{G}\boldsymbol{\beta}$. You predict selection will move the population northeast; it moves east. The discrepancy is systematic, not random.
\item Fourth, environmental sensitivity increases unexpectedly. G$\times$E variance is growing as a fraction of total variance. The gradient is rotating relative to your genetic structure.
\item Fifth, realized heritability drops. You select hard but response is weak. The trait looks less heritable than it did before.
\item Any one of these could have other explanations. But if you're seeing several together, you should investigate. You might be leaving an additive channel.
\item The good news is all of these are measurable from standard breeding or evolutionary data. You don't need new experiments---you need to look at your existing data in a new way.
}

% =============================================================================
% PART VI: SYNTHESIS (Slides 29–31)
% =============================================================================
\section{Part VI: Synthesis}

% -----------------------------------------------------------------------------
% SLIDE 29: Resolving the Paradox
% -----------------------------------------------------------------------------
\begin{frame}{Resolving the Paradox}
\frametitle{Resolving the Paradox}

\textbf{Why Additive Models Work on Non-Additive Architectures}

\begin{enumerate}
    \item \textbf{Local linearity:} Any smooth function is locally linear
    \item \textbf{Scaling:} $\Vquad \propto \sigma^4$ is negligible at small variance
    \item \textbf{Self-correction:} Selection compresses variance, raising $\Ag$
    \item \textbf{Transient curves:} Escape from additivity is brief
\end{enumerate}

\vspace{0.5cm}
\begin{alertblock}{}
Additivity is not a property of the genome. It's a property of how much of the landscape the population samples.
\end{alertblock}

\end{frame}

\note[itemize]{
\item \textbf{[2 minutes]} Let me now bring everything together and resolve the paradox we started with.
\item We asked: how can additive models work when genomes are non-additive?
\item The answer rests on four pillars.
\item First, \textbf{local linearity}. Any smooth function is locally linear. The fitness landscape can be globally curved, but near any point, the tangent plane is a good approximation. This is just calculus.
\item Second, \textbf{scaling}. The quadratic variance scales as $\sigma^4$, not $\sigma^2$. When genetic variance is small, $\sigma^4$ is very small. The curvature term contributes almost nothing. The world looks linear.
\item Third, \textbf{self-correction}. Selection compresses variance. And when variance shrinks, the quadratic term shrinks faster than the linear term. Selection naturally pushes populations into additive channels. Additivity is self-reinforcing.
\item Fourth, \textbf{transient curves}. Even when populations do escape---through introgression, environmental shifts, or radiation events---the escape is brief. They negotiate the curve and settle into a new channel.
\item The resolution is in that boxed statement. \textbf{Additivity is not a property of the genome. It's a property of how much of the landscape the population samples.} Narrow sampling gives local linearity. Wide sampling exposes curvature.
\item Epistatic architectures produce additive-looking dynamics because selection keeps populations in regimes where the linear approximation works.
}

% -----------------------------------------------------------------------------
% SLIDE 30: Key Take-Home Messages
% -----------------------------------------------------------------------------
\begin{frame}{Key Take-Home Messages}
\frametitle{Key Take-Home Messages}

\textbf{What to Remember}

\begin{enumerate}
    \item $\Ag = \Vlin/(\Vlin + \Vquad)$ measures ``local additivity''
    \item $\Ag = R^2$ of the linear predictor for log-fitness
    \item Fisher's theorem, the breeder's equation, and GBLUP work when $\Ag$ is high
    \item Selection naturally pushes populations into additive channels
    \item Wide crosses, new environments, and radiations trigger escape
    \item $\Ag$ can be estimated and used as a breeding/evolution diagnostic
\end{enumerate}

\end{frame}

\note[itemize]{
\item \textbf{[1.5 minutes]} If you remember six things from this talk, let it be these.
\item One: the additivity index $\Ag$ measures how linear the world looks to your population. It's the fraction of fitness variance captured by the linear term.
\item Two: under Gaussian assumptions, $\Ag$ equals the $R^2$ of the linear predictor. It's directly interpretable as prediction accuracy.
\item Three: classical results---Fisher's theorem, the breeder's equation, GBLUP accuracy---are local properties. They work when $\Ag$ is high. They fail when it's low.
\item Four: selection is self-correcting. It compresses variance and raises $\Ag$. Populations naturally find additive channels.
\item Five: escape happens through variance expansion or gradient rotation. Wide crosses, novel environments, adaptive radiations---these push populations out of channels.
\item Six: $\Ag$ can be estimated from real data and used as a diagnostic. It tells you whether your bottleneck is data quality or model specification.
\item The framework unifies landscape theory with practical prediction. It explains why additive models work, when they fail, and what to do about it.
}

% -----------------------------------------------------------------------------
% SLIDE 31: Future Directions & Questions
% -----------------------------------------------------------------------------
\begin{frame}{Future Directions \& Questions}
\frametitle{Future Directions \& Questions}

\textbf{Where Do We Go From Here?}

\begin{itemize}
    \item Empirical estimation of $\Ag$ across taxa/traits
    \item Non-Gaussian extensions (selection truncation, skewness)
    \item Dynamic $\mathbf{G}$: tracking channels through time
    \item Breeding applications: explicit $\Ag$ optimisation?
    \item Connections to fitness landscape theory (DFE, rugged landscapes)
\end{itemize}

\vspace{1cm}
\begin{center}
\Large \textbf{Thank you!} \\[0.5cm]
\normalsize Questions?
\end{center}

\end{frame}

\note[itemize]{
\item \textbf{[1 minute]} Let me close with some open questions.
\item First, \textbf{empirical estimation}. We need to measure $\Ag$ across different taxa, traits, and breeding programs. Is it typically high? When is it low? What predicts it?
\item Second, \textbf{non-Gaussian extensions}. Selection truncation and other forces create skewed distributions. How does that affect the orthogonality between linear and quadratic terms? How does it change the $R^2$ interpretation?
\item Third, \textbf{dynamic $\mathbf{G}$ matrices}. How do we track populations as they move through channels over time? Can we predict when channels will narrow or bend?
\item Fourth, \textbf{breeding optimisation}. Should we design programs to maintain high $\Ag$? Or is there value in deliberately pushing into low-$\Ag$ regions to explore new variation?
\item Fifth, \textbf{connections to fitness landscape theory}. How does $\Ag$ relate to the distribution of fitness effects? To measures of ruggedness? To evolutionary accessibility?
\item These are questions for the next papers. For now, I hope I've convinced you that the additivity puzzle has a geometric solution, and that the solution is useful.
\item Thank you. I'm happy to take questions.
}

% =============================================================================
% SUPPLEMENTARY SLIDES (for Q&A)
% =============================================================================
\appendix
\section{Supplementary Slides}

% -----------------------------------------------------------------------------
% SLIDE S1: Mathematical Details – Variance of Quadratic Forms
% -----------------------------------------------------------------------------
\begin{frame}{Mathematical Details -- Variance of Quadratic Forms}
\frametitle{Mathematical Details -- Variance of Quadratic Forms}

\textbf{For the Technically Curious}

For $\mathbf{x} \sim N(\mathbf{0}, \boldsymbol{\Sigma})$ and symmetric $\mathbf{A}$:
\[
\var(\mathbf{x}^\top \mathbf{A} \mathbf{x}) = 2\,\tr[(\mathbf{A}\boldsymbol{\Sigma})^2]
\]

Therefore:
\[
\Vquad = \var\!\left(\tfrac{1}{2}\Delta\mathbf{g}^\top \mathbf{H} \Delta\mathbf{g}\right) = \tfrac{1}{2}\tr[(\mathbf{H}\boldsymbol{\Sigma}_g)^2]
\]

This is why $\Vquad \propto \sigma^4$ (each $\boldsymbol{\Sigma}$ contributes $\sigma^2$).

\end{frame}

\note[itemize]{
\item \textbf{[Use if asked about the $\sigma^4$ scaling]}
\item This is a standard result for quadratic forms under normality. If $\mathbf{x}$ is multivariate normal with covariance $\boldsymbol{\Sigma}$, and $\mathbf{A}$ is symmetric, the variance of the quadratic form $\mathbf{x}^\top \mathbf{A} \mathbf{x}$ is twice the trace of $(\mathbf{A}\boldsymbol{\Sigma})^2$.
\item Each factor of $\boldsymbol{\Sigma}$ contributes one power of $\sigma^2$. The trace of the square means you multiply two factors together. Hence $\sigma^4$.
\item The half in front comes from the $\frac{1}{2}$ in the Taylor expansion of the quadratic term. When you take the variance, the half gets squared to give a quarter, but the formula has a factor of 2, so you end up with $\frac{1}{2}$.
\item This is why the scaling insight works. It's not an approximation---it's an exact result for Gaussian distributions.
}

% -----------------------------------------------------------------------------
% SLIDE S2: Orthogonality of Linear and Quadratic Terms
% -----------------------------------------------------------------------------
\begin{frame}{Orthogonality of Linear and Quadratic Terms}
\frametitle{Orthogonality of Linear and Quadratic Terms}

\textbf{Why $\cov(L, Q) = 0$ Under Gaussianity}

For jointly normal $\mathbf{x}$ with mean zero:
\[
\cov(\mathbf{a}^\top \mathbf{x}, \mathbf{x}^\top \mathbf{B} \mathbf{x}) = 0
\]

This is a property of Gaussian distributions only. It guarantees that:
\[
\var(\tilde{\ell}) = \var(L) + \var(Q) = \Vlin + \Vquad
\]

Non-Gaussianity (e.g., from selection truncation) breaks this.

\end{frame}

\note[itemize]{
\item \textbf{[Use if asked about the $R^2$ interpretation or Gaussian assumptions]}
\item This orthogonality is crucial for the $R^2$ interpretation. Under Gaussianity, a linear function and a quadratic function of the same normal vector are uncorrelated.
\item Intuitively, the linear function is odd in each component---it changes sign when you flip the sign of $x_i$. The quadratic function is even. Their product integrates to zero over a symmetric distribution.
\item This means the variance decomposition is additive: $\var(L + Q) = \var(L) + \var(Q)$. No covariance term.
\item But this breaks under non-Gaussianity. If you've truncated the distribution by selection, the linear and quadratic terms can be correlated. The $R^2$ interpretation becomes approximate rather than exact.
\item In practice, this is probably a second-order effect for most breeding populations. But it's a limitation to be aware of.
}

% -----------------------------------------------------------------------------
% SLIDE S3: The 2D Example – Alignment Effects
% -----------------------------------------------------------------------------
\begin{frame}{The 2D Example -- Alignment Effects}
\frametitle{The 2D Example -- Alignment Effects}

\textbf{Orientation of Genetic Variance Relative to Curvature}

In 2D, $\Ag$ depends on how the eigenvectors of $\boldsymbol{\Sigma}_g$ align with those of $\mathbf{H}$.

\begin{itemize}
    \item \textbf{Aligned:} curvature is sampled maximally $\to$ low $\Ag$
    \item \textbf{Orthogonal:} variance lies in ``flat'' directions $\to$ high $\Ag$
\end{itemize}

\vspace{0.5cm}
This is why $\mathbf{G}$ rotation matters: the same total variance can give very different $\Ag$ values.

\end{frame}

\note[itemize]{
\item \textbf{[Use if asked about why $\mathbf{G}$ matrix rotation matters]}
\item In higher dimensions, $\Ag$ depends not just on the total variance but on how that variance aligns with curvature.
\item Imagine an elongated ellipse of genetic variation and a fitness surface that curves strongly along one axis but not another.
\item If the ellipse aligns with the curved direction, the population samples a lot of curvature. $\Vquad$ is large. $\Ag$ is low.
\item If the ellipse is perpendicular---variance along the flat direction---the population avoids the curvature. $\Vquad$ is small. $\Ag$ is high.
\item Same total variance. Very different $\Ag$.
\item This is why $\mathbf{G}$ matrix rotation matters for additivity, not just for evolvability. The same genetic variance can look additive or non-additive depending on its orientation relative to the landscape.
}

% -----------------------------------------------------------------------------
% SLIDE S4: Connection to Bulmer Effect
% -----------------------------------------------------------------------------
\begin{frame}{Connection to Bulmer Effect}
\frametitle{Connection to Bulmer Effect}

\textbf{Selection Reduces Genetic Variance}

The Bulmer effect: directional selection reduces $V_A$ within a generation.

In our framework:
\begin{itemize}
    \item Selection compresses $\boldsymbol{\Sigma}_g$ along the gradient
    \item $\Vquad$ drops faster than $\Vlin$
    \item $\Ag$ increases after each selection cycle
\end{itemize}

\vspace{0.5cm}
The Bulmer effect is part of why populations stay in additive channels.

\end{frame}

\note[itemize]{
\item \textbf{[Use if asked about dynamics of variance compression]}
\item The Bulmer effect is the within-generation reduction in genetic variance due to selection. Before selection, you have full variance. After selection, the selected parents have reduced variance because extreme genotypes were preferentially chosen.
\item In our framework, this is the mechanism by which selection compresses $\boldsymbol{\Sigma}_g$. Each round of selection squeezes the ellipse of genetic variation, particularly along the direction of the selection gradient.
\item Because $\Vquad$ scales as $\sigma^4$, this compression has an outsized effect on the quadratic term. Each cycle, the quadratic contribution shrinks faster than the linear contribution.
\item This is the self-correcting dynamic in action. The Bulmer effect, often seen as a nuisance that limits long-term response, is actually helping maintain additivity. It pushes populations deeper into additive channels.
\item Recombination and mutation eventually restore variance, but over relevant timescales, selection tends to win. That's why additive models work.
}

% =============================================================================
\end{document}
% =============================================================================
